\section{生成式视觉语言模型设计}

本章在语义分类任务已验证 EEG 可提供有效辅助信息的基础上，将经 VTF 融合后的 EEG-enhanced vision tokens 接入生成式视觉语言模型，构建从退化图像与脑电信号到自然语言描述的整体生成流程。

\subsection{生成式描述任务的设计目标}

本任务要求模型根据输入的退化图像及其对应的 EEG trial，输出一段对该图像内容进行描述的自然语言 caption。与第 5 章所处理的语义分类不同，此处模型不再输出预定义的类别标签，而是生成自由形式的文本描述。

这里需要澄清的是，本设计的出发点并非``通过脑电信号直接解读任意语义内容''。本课程设计在约 40 个类别、12000 条 trial 的数据规模下，将 EEG 定位为一种视觉语义增强信号——利用 EEG 所携带的、与图像内容相关的额外神经表征信息，使生成式视觉语言模型在面对退化输入时能够产生比仅依赖视觉信息更为准确的描述。

\subsection{轻量级 GRU 图像描述基线}

在引入大规模预训练生成模型之前，本设计先搭建了一个轻量级的 GRU 自回归 caption 解码器作为基线方案。该解码器以 VTF 输出的 enhanced vision tokens 经平均池化后的向量作为初始隐状态，采用单层 GRU 逐 token 解码文本序列。模型总参数量约 2.8M，在 BLIP 生成的伪 caption 上完成训练。

GRU 基线模型能够稳定地产出 free-form caption，其优势在于训练开销低、无需依赖大规模预训练模型、输出文本长度始终自然可控。然而，GRU 解码器缺乏足够的语言先验知识，生成结果的语法流畅度和概念准确性均显著弱于预训练视觉语言模型。GRU 基线的主要意义在于验证了 EEG-enhanced vision tokens 具备驱动自回归 caption 生成的能力。

\subsection{预训练生成模型方案比较}

为确定最优的生成方案，本设计系统对比了五种技术路线，结果汇总于表~\ref{tab:route-comparison}。实验过程中为方便记录，对各方案使用了简明代号。

\begin{table}[htbp]
  \centering
  \caption{生成式视觉语言模型方案对比}
  \label{tab:route-comparison}
  \begin{tabularx}{\linewidth}{l l X}
    \toprule
    方案代号 & 输入与桥接方式 & 主要发现 \\
    \midrule
    方案一 & enhanced tokens 池化后投影为 inputs\_embeds，输入 Qwen2.5 + LoRA & 纯文本 LLM 缺乏视觉-语言对齐能力，有效输出率低 \\
    方案二 & Q-Former/Perceiver bridge 桥接视觉与语言 & 在小样本条件下效果有限，需要更大的训练数据 \\
    方案三 & LLaVA-style MLP projector & 与 Qwen2-VL adapter 的工程兼容性不佳 \\
    方案四 & BLIP-2 style Q-Former & 清洁图像上表现尚可，退化条件下不稳定 \\
    方案五 & enhanced tokens + Qwen2-VL adapter/prefix + LoRA & 最终采用的生成方案 \\
    \bottomrule
  \end{tabularx}
\end{table}

综合各项指标来看，方案五（基于 Qwen2-VL adapter）在有效输出率、caption class-hit 以及 real-control gap 三个维度上均领先于其他四种方案。Qwen2-VL 作为原生视觉语言模型，其视觉编码模块在预训练阶段已与语言模型实现充分的图文对齐，因而在小样本微调场景下能够更有效地利用 EEG-enhanced visual tokens。

\subsection{基于 Qwen2-VL 的最终生成方案}

本设计最终采用的生成方案架构如图~\ref{fig:route5} 所示。

\placeholderfigure{基于 Qwen2-VL 的生成式视觉语言模型架构图}{route5-qwenvl-placeholder}{5.0cm}

其核心处理流程如下：(1) 退化图像经 CLIP ViT 提取 visual tokens，再与 EEG tokens 通过 VTF 的 cross-attention 机制融合，获得增强视觉 token；(2) 增强视觉 token 先经本设计实现的轻量 prefix adapter 投影到 Qwen2-VL 可接受的嵌入空间，随后与语义提示文本一同送入生成式解码器；(3) 在输入序列中附加语义提示信息（A2 模型预测的最可能类别）和退化类型标签，以文本 token 形式拼接到生成 prompt 中；(4) Qwen2-VL-2B-Instruct 解码器在 teacher-forcing 模式下进行自回归训练，测试阶段则以 ``corruption\_type:'' 为前缀触发文本生成。

\subsection{LoRA 微调与训练配置}

生成式模型采用 LoRA 方法，在 Qwen2-VL 的 Query 和 Value 投影矩阵上施加低秩适配，其中秩 $r=8$，缩放因子 $\alpha=16$。可训练参数量约为 1.3M，占模型总体参数的比例极小。训练过程中冻结所有原始模型权重。本设计还额外对比了 $r=16$ 的配置效果。

在训练目标方面，本设计测试了两种策略：(1) 以类别名称为优化目标——直接针对类别正确性进行训练；(2) 以类别名称配合经过滤的 BLIP 伪 caption 为优化目标——生成的文本更自然流畅，但 BLIP 伪 caption 中可能存在的噪声会对类别准确性带来负面影响。实验结果表明，类别目标更有利于提升 caption 的类别命中率，而类别配合 BLIP 伪 caption 的目标则更有利于改善文本自然度和提高有效输出率。本设计分别报告这两类目标的实验效果，最终以 Qwen2-VL LoRA $r=8$ 配合多候选重排序作为主要的生成式结果呈现。

\subsection{多候选重排序策略}

推理阶段，本设计为每个输入生成 $N=5$ 个候选 caption（通过调节温度参数、top-p 等解码策略实现多样化输出），随后依据候选 caption 的有效性（非空、非 URL 或代码片段）、重复率和长度进行初步筛选。筛选后，计算各候选与 A2/VTF 预测的 top-k 语义概念或 CLIP 文本原型之间的语义相似度，选取一致性最高的候选项作为最终输出。整个过程不依赖测试集真实类别标签，caption 类别命中率仅用于最终阶段的评估。多候选重排序机制在不改变模型参数的前提下，显著提升了有效输出率以及定性展示的效果。
